\providecommand{\todo}[1]{\textcolor{red}{\textbf{[TODO: #1]}}}
\clearpage
\phantomsection
\section*{Chương 8. Chuẩn hóa $C[\tau]$ và octave cost}
\addcontentsline{toc}{section}{Chương 8. Chuẩn hóa $C[\tau]$ và octave cost}

\begin{learningobj}
\begin{itemize}[leftmargin=1.5em, itemsep=2pt]
    \item Thấu hiểu bản chất vật lý và nguồn gốc toán học của hiện tượng suy giảm hình học trong hàm tự tương quan ngắn hạn: khi độ trễ $\tau$ tăng, diện tích chồng lấn của cửa sổ co lại khiến năng lượng tự tương quan của cửa sổ $R_w[\tau]$ giảm dần từ $1{,}0$ về $0{,}468$ tại $\tau = 214$; phân tích sự thật phũ rằng nếu không chia bù cho $R_w[\tau]$, các đỉnh ở độ trễ lớn bị dìm xuống khiến thuật toán thiên vị chọn độ trễ nhỏ (tần số cao), dẫn đến lỗi nhảy quãng tám lên trên (octave doubling).
    \item Khảo sát nguyên lý chi phí quãng tám (Octave Cost) trong mã nguồn Praat 7.0.03 (\path{fon/Sound_to_Pitch.cpp}, dòng 158--164, 225--240 và \path{fon/Pitch.cpp}, dòng 554--560); giải mã lý do Paul Boersma trừ một lượng phạt tỷ lệ với logarit chu kỳ nhằm ưu tiên tần số cao và phá vỡ thế cân bằng nguy hiểm giữa chu kỳ cơ bản $\tau_0$ và chu kỳ bội $2\tau_0$ (octave halving); chứng minh bằng giải tích sự tương đương giữa công thức phân loại ứng viên theo $F_{\text{floor}}$ và hàm chi phí cục bộ theo $F_{\text{ceiling}}$ (lệch nhau đúng một hằng số cộng $0{,}030000$).
    \item Phân tích toàn diện ba phương án thiết kế phần cứng: chứng minh giải pháp chia trực tiếp $R[\tau] / R_w[\tau]$ tiêu tốn tài nguyên quá mức; vạch rõ sai số của giải pháp gộp chi phí vào ROM nhân (phiên bản v1) làm nghiêng đỉnh parabol và dịch tần số từ $0{,}07$ đến $0{,}23\text{ cent}$, đồng thời chỉ rõ công thức của Argon không thể tạo ra hệ số $17070$; bảo vệ quyết định chốt kiến trúc trừ đúng như Praat qua bộ bảng tra ROM tích hợp $256 \times 32$ trong đúng 1 khối M10K (\path{golden_v2/rom/ck_rom.{hex,mif}}).
    \item Dẫn xuất bước chuẩn hóa động $\text{sh} = \max(0, \text{bitlen}(R_0) - 26)$ đưa toán hạng $R_n[\tau]$ về $27\text{ bit}$ có dấu, chứng minh sai số tương đối của phép dịch số học này không vượt quá $2^{-26} \approx 1{,}49 \times 10^{-8} \le 1{,}5 \times 10^{-8}$; giải thích cơ chế silicon giúp hai phép nhân $R_n \times C$ ($27 \times 17$) và $R_{0n} \times K$ ($27 \times 16$) nằm vừa khít trong đúng 2 khối Variable Precision DSP $27 \times 27$ của Cyclone V.
    \item Thiết kế và hiện thực hóa hai tuyến phần cứng: Tuyến (A) cấu hình khối IP \texttt{ROM: 1-PORT} (\emph{Embedded Memory User Guide}, UG-01068) và khối nhân \texttt{LPM\_MULT} (\emph{Integer Arithmetic IP Cores User Guide}, UG-01069); Tuyến (B) module HDL tự viết \texttt{ch08_ck_rom.v} và khối chuẩn hóa tính điểm đường ống 3 nhịp \texttt{ch08_norm_score.v}.
    \item Thực thi quy trình kiểm chứng chu-kỳ-đối-chu-kỳ giữa mô hình song sinh Python (\texttt{ch08_twin.py}) và mô phỏng Verilog (\texttt{tb_ch08_norm.v}) trên $1.214\text{ vector}$ kích thích bao gồm khung sạch, khung ồn SNR 5 dB, khung cực hạn stress full-scale ($R_0 \approx 2^{42}$), khung biên độ nhỏ và khung im lặng; chứng minh $U[\tau] / R_{0n}$ tái lập trung thực hàm $r[\tau]$ của Praat 7.0.03 với sai số trung bình chỉ $2{,}12 \times 10^{-3}$ trên tín hiệu thực tế.
\end{itemize}
\end{learningobj}

\subsection*{8.1. Bản chất của việc chuẩn hóa cửa sổ và tại sao hàm tự tương quan suy giảm theo độ trễ}
\addcontentsline{toc}{subsection}{8.1. Bản chất của việc chuẩn hóa cửa sổ và tại sao hàm tự tương quan suy giảm theo độ trễ}

\subsubsection*{Vấn đề thực tế: Sự suy giảm hình học nhân tạo và nguy cơ nhảy quãng tám}
Hãy tưởng tượng một tình huống cụ thể trong xử lý tiếng nói: người nói phát âm một nguyên âm trầm ở tần số cơ bản $F_0 = 100\text{ Hz}$. Tại tần số lấy mẫu $f_s = 16.000\text{ Hz}$, chu kỳ dao động cơ bản của dây thanh âm là $\tau_0 = 16.000 / 100 = 160\text{ mẫu}$. Khi tín hiệu âm thanh này được nhân với cửa sổ Hanning chiều dài hữu hạn $N = 638\text{ mẫu}$ rồi đưa vào khối MAC ở Chương 7 để tính hàm tự tương quan $R[\tau]$, điều gì sẽ xảy ra tại các độ trễ $\tau$?

Theo trực giác thuần túy, vì dạng sóng lặp lại tuần hoàn sau mỗi $160\text{ mẫu}$, ta kỳ vọng tích vô hướng giữa dạng sóng với bản sao trễ $160\text{ mẫu}$ phải có giá trị tương đương với năng lượng gốc tại $\tau = 0$. Nhưng thực tế phần cứng lại đưa ra một kết quả gây bối rối: giá trị $R[160]$ thu được chỉ đạt khoảng $65{,}4\%$ so với $R[0]$! Thậm chí, nếu tín hiệu có họa ba bậc hai mạnh ở $200\text{ Hz}$ (ứng với độ trễ $\tau_1 = 80\text{ mẫu}$), giá trị $R[80]$ lại đạt tới $79{,}1\%$ giá trị đỉnh. 

Nếu thuật toán chọn đỉnh cực đại trực tiếp từ dãy số $R[\tau]$ chưa chuẩn hóa này, đỉnh tại $\tau = 80\text{ mẫu}$ sẽ chiến thắng đỉnh tại $\tau = 160\text{ mẫu}$. Kết quả là hệ thống sẽ báo cao độ là $200\text{ Hz}$ thay vì $100\text{ Hz}$! Hiện tượng nhận diện sai chu kỳ cơ bản thành một nửa chu kỳ thật (tức tần số bị nhân đôi) được gọi là \textbf{lỗi nhảy quãng tám lên trên} (octave doubling).

Nguyên nhân của sự sai lệch này không nằm ở bản chất giọng nói, mà nằm ở \textbf{hình học của việc cắt khung hữu hạn}. Hàm tự tương quan của tín hiệu đã nhân cửa sổ $x_w[n] = x[n] \cdot w[n]$ được xác định bởi:
\begin{equation}
R_{xw}[\tau] = \sum_{n=0}^{N-1-\tau} x_w[n] \cdot x_w[n+\tau] = \sum_{n=0}^{N-1-\tau} x[n] x[n+\tau] \cdot w[n] w[n+\tau]
\label{eq:ch08_rxw_prod}
\end{equation}

Khi tín hiệu $x[n]$ là một quá trình dừng có tính tuần hoàn mạnh, tích $x[n] x[n+\tau]$ dao động nhanh xung quanh giá trị trung bình $R_x[\tau]$, trong khi tích hai mẫu cửa sổ $w[n] w[n+\tau]$ biến thiên chậm rãi như một đường bao hình vòm. Áp dụng nguyên lý tách biến (separation of scales), tổng tích lũy~(\ref{eq:ch08_rxw_prod}) có thể xấp xỉ chặt chẽ bằng tích của hàm tự tương quan tín hiệu gốc với hàm tự tương quan của chính cửa sổ:
\begin{equation}
R_{xw}[\tau] \approx R_x[\tau] \cdot \frac{R_w[\tau]}{R_w[0]}
\label{eq:ch08_separation}
\end{equation}
trong đó hàm tự tương quan của cửa sổ Hanning được định nghĩa là:
\begin{equation}
R_w[\tau] = \sum_{n=0}^{N-1-\tau} w[n] \cdot w[n+\tau]
\label{eq:ch08_rw_def}
\end{equation}

\subsubsection*{Dẫn xuất toán học: Sự co ngắn của diện tích chồng lấn}
Hãy quan sát cận trên của tổng trong phương trình~(\ref{eq:ch08_rw_def}): số lượng số hạng tham gia vào phép tính tự tương quan là đúng $N - \tau\text{ mẫu}$. 
\begin{itemize}[leftmargin=1.5em, itemsep=2pt]
    \item Khi $\tau = 0$: Hai cửa sổ Hanning trùng khít hoàn toàn trên toàn bộ $N = 638\text{ mẫu}$. Đại lượng $R_w[0] = \sum_{n=0}^{637} w^2[n]$ đạt giá trị cực đại. Với cửa sổ Hanning $N+1$ đã lượng tử hóa Q1.15 ở Chương 6, ta có $R_w[0] = 257.295.365.938$ (nhật ký \path{verify_v3/ch08_norm_math.log}).
    \item Khi $\tau$ tăng dần từ $26$ đến $214$: Miền giao nhau giữa cửa sổ gốc và cửa sổ dịch trễ bị thu hẹp dần từ $612\text{ mẫu}$ xuống chỉ còn $424\text{ mẫu}$. Đồng thời, các mẫu có trọng số lớn nhất ở vùng đỉnh trung tâm ($w[318] = w[319] = 32767$) bị lệch pha nhau, khiến tích $w[n] w[n+\tau]$ suy giảm nghiêm trọng.
\end{itemize}

Tại cận trên của dải quét phần cứng $\tau = 214$, giá trị $R_w[214]$ chỉ còn:
\begin{equation}
\frac{R_w[214]}{R_w[0]} = \frac{120.407.728.314}{257.295.365.938} \approx \mathbf{0{,}46797} \approx 46{,}8\%
\end{equation}
Nói cách khác, một thành phần sóng có chu kỳ dài $\tau = 214\text{ mẫu}$ ($F_0 \approx 74{,}77\text{ Hz}$) bị dìm biên độ tự nhiên mất hơn $53{,}2\%$ năng lượng chỉ vì hiệu ứng mép cửa sổ! 

\begin{firstprin}{Sự cần thiết tuyệt đối của việc chia bù hàm tự tương quan cửa sổ}
Để khôi phục lại hàm tự tương quan khách quan của tín hiệu $R_x[\tau]$, ta bắt buộc phải triệt tiêu ảnh hưởng hình học của cửa sổ bằng cách chia giá trị tự tương quan đo được $R_{xw}[\tau]$ cho hàm tự tương quan chuẩn hóa của cửa sổ $R_w[\tau] / R_w[0]$:
\begin{equation}
r[\tau] = \frac{R_{xw}[\tau]}{R_{xw}[0] \cdot \left( \frac{R_w[\tau]}{R_w[0]} \right)} = \frac{R_{xw}[\tau]}{R_{xw}[0]} \cdot \frac{R_w[0]}{R_w[\tau]}
\label{eq:ch08_norm_exact}
\end{equation}
Đại lượng $r[\tau]$ chính là hàm tự tương quan chuẩn hóa không thứ nguyên (normalized autocorrelation), có giá trị lý thuyết $r[0] = 1{,}0$ và $|r[\tau]| \le 1{,}0$ trên toàn bộ dải trễ.
\end{firstprin}

Trong phần mềm Praat 7.0.03, Paul Boersma đã hiện thực hóa nguyên lý này tại tệp \path{fon/Sound_to_Pitch.cpp}, dòng 158--164:
\begin{lstlisting}[language=C++, numbers=none]
/* Normalize the autocorrelation to the value with zero lag,
   and divide it by the normalized autocorrelation of the window. */
r [0] = 1.0;
for (integer i = 1; i <= brent_ixmax; i ++)
    r [i] = ac [i + 1] / (ac [1] * windowR [i + 1]);
\end{lstlisting}
Trong đó mảng \texttt{windowR} lưu trữ sẵn giá trị $R_w[\tau] / R_w[0]$ được tính toán trước qua biến đổi FFT (dòng 427--442), và \texttt{ac[1]} tương ứng với giá trị năng lượng khung $R_{xw}[0]$.

\subsubsection*{Tính toán thủ công thu nhỏ: Cửa sổ Hanning 8 mẫu}
Để thấy rõ cơ chế bù trừ này bằng trực giác số học, ta hãy khảo sát một ví dụ tính toán thu nhỏ với cửa sổ Hanning dạng Praat $N+1$ gồm $N = 8\text{ mẫu}$. Công thức giải tích là:
\begin{equation}
w[i] = 0{,}5 - 0{,}5 \cos\left(\frac{2\pi i}{N+1}\right) = 0{,}5 - 0{,}5 \cos\left(\frac{2\pi i}{9}\right), \quad i = 1, 2, \dots, 8
\end{equation}
Các giá trị mẫu của cửa sổ $w[i]$ đối xứng qua tâm:
\begin{equation}
w = [0{,}116978;\; 0{,}413176;\; 0{,}750000;\; 0{,}969846;\; 0{,}969846;\; 0{,}750000;\; 0{,}413176;\; 0{,}116978]
\end{equation}

Năng lượng gốc tại $\tau = 0$ là tổng bình phương 8 mẫu:
\begin{equation}
R_w[0] = \sum_{i=1}^8 w^2[i] = 2 \times (0{,}116978^2 + 0{,}413176^2 + 0{,}750000^2 + 0{,}969846^2) = \mathbf{3{,}375000}
\end{equation}

Tại các độ trễ $\tau = 1, 2, 3$, tích vô hướng giữa cửa sổ và bản sao dịch chuyển của nó suy giảm nhanh chóng:
\begin{itemize}[leftmargin=1.5em, itemsep=2pt]
    \item Tại $\tau = 1$: $R_w[1] = \sum_{i=1}^7 w[i] w[i+1] \approx 3{,}111800 \implies R_w[1]/R_w[0] \approx 0{,}922015$.
    \item Tại $\tau = 2$: $R_w[2] = \sum_{i=1}^6 w[i] w[i+2] \approx 2{,}431670 \implies R_w[2]/R_w[0] \approx 0{,}720495$.
    \item Tại $\tau = 3$: $R_w[3] = \sum_{i=1}^5 w[i] w[i+3] \approx 1{,}590835 \implies R_w[3]/R_w[0] \approx 0{,}471359$.
\end{itemize}

Hệ số khuếch đại bù suy hao lý tưởng là nghịch đảo $C_{\text{ideal}}[\tau] = R_w[0] / R_w[\tau]$:
\begin{equation}
C_{\text{ideal}}[0] = 1{,}000000; \quad C_{\text{ideal}}[1] = 1{,}084581; \quad C_{\text{ideal}}[2] = 1{,}387935; \quad C_{\text{ideal}}[3] = 2{,}121527
\end{equation}

Bây giờ, ta cấp vào hệ thống một tín hiệu số nguyên tuần hoàn có chu kỳ cơ bản đúng $\tau_0 = 3\text{ mẫu}$:
\begin{equation}
x = [2,\; -1,\; -2,\; 2,\; -1,\; -2,\; 2,\; -1]
\end{equation}
Nhân tín hiệu $x$ với cửa sổ $w$, ta thu được tín hiệu có cửa sổ $x_w[i] = x[i] \cdot w[i]$:
\begin{equation}
x_w = [0{,}233956;\; -0{,}413176;\; -1{,}500000;\; 1{,}939693;\; -0{,}969846;\; -1{,}500000;\; 0{,}826352;\; -0{,}116978]
\end{equation}

Tính hàm tự tương quan ngắn hạn chưa chuẩn hóa $R_{xw}[\tau]$:
\begin{itemize}[leftmargin=1.5em, itemsep=2pt]
    \item Tại $\tau = 0$: $R_{xw}[0] = \sum x_w^2[i] = 10{,}125000$.
    \item Tại $\tau = 3$: $R_{xw}[3] = \sum_{i=1}^5 x_w[i] x_w[i+3] = 4{,}820838$.
\end{itemize}

Quan sát tỉ số chưa chuẩn hóa:
\begin{equation}
\frac{R_{xw}[3]}{R_{xw}[0]} = \frac{4{,}820838}{10{,}125000} \approx \mathbf{0{,}476132}
\end{equation}
Dù tín hiệu $x$ lặp lại tuần hoàn hoàn hảo sau mỗi 3 mẫu, tỉ số tương quan đo được bị tụt từ $1{,}0$ xuống dưới $0{,}48$! Nếu một thành phần nhiễu tần số cao tại $\tau = 1$ có tỉ số tương quan giả mạo cỡ $0{,}5$, thuật toán sẽ lập tức chọn nhầm $\tau = 1$.

Tuy nhiên, khi ta áp dụng phép chuẩn hóa cửa sổ theo phương trình~(\ref{eq:ch08_norm_exact}):
\begin{equation}
r[3] = \frac{R_{xw}[3]}{R_{xw}[0]} \times C_{\text{ideal}}[3] = 0{,}476132 \times 2{,}121527 = \mathbf{1{,}010127} \approx 1{,}0
\end{equation}
(Giá trị hơi vượt $1{,}0$ do hiệu ứng lượng tử hóa biên trên khung siêu ngắn 8 mẫu; trên khung 638 mẫu của Praat, giá trị này hội tụ chặt chẽ về $1{,}0$). Nhờ nhân bù với $C_{\text{ideal}}[3] = 2{,}1215$, độ mạnh của đỉnh tại chu kỳ thật $\tau_0 = 3$ đã được khôi phục trọn vẹn!

\begin{figure}[htbp]
\centering
\resizebox{0.95\textwidth}{!}{% chuong_v3/fig_ch08_window_effect.tex
% Minh họa hiệu ứng suy giảm hình học của hàm tự tương quan cửa sổ Hanning Rw[tau]
% và cơ chế bù suy hao qua hệ số C[tau] = round(2^14 * Rw[0] / Rw[tau])
\begin{tikzpicture}[x=0.06cm, y=4.5cm, >=Stealth]
    % Lưới nền mờ
    \draw[very thin, gray!20, step=1] (0,-0.5) grid[xstep=30, ystep=0.2] (225,1.2);

    % Trục tọa độ
    \draw[->, thick, black!70] (0,0) -- (235,0) node[right, font=\footnotesize] {Độ trễ $\tau$ (mẫu)};
    \draw[->, thick, black!70] (0,-0.55) -- (0,1.25) node[above, font=\footnotesize] {Biên độ chuẩn hóa};

    % Vạch chia trục hoành
    \foreach \x in {0, 26, 60, 100, 140, 180, 214} {
        \draw[thick, black!60] (\x,0.02) -- (\x,-0.02) node[below, font=\scriptsize, fill=white, inner sep=1pt] {\x};
    }
    % Vạch chia trục tung
    \foreach \y in {-0.4, -0.2, 0.0, 0.2, 0.4, 0.6, 0.8, 1.0} {
        \draw[thick, black!60] (2,\y) -- (-2,\y) node[left, font=\scriptsize] {\y};
    }

    % Vùng quét trễ tau in [26, 214]
    \fill[blue!5, opacity=0.7] (26,-0.52) rectangle (214,1.16);
    \node[blue!70!black, font=\scriptsize, fill=white, inner sep=2.5pt, rounded corners=2pt, draw=blue!30] at (120,1.16) 
        {Dải quét trễ phần cứng: $\tau \in [26, 214]$ ($F_0 \in [75, 600]\text{ Hz}$)};

    % Đường cong 1: Tự tương quan của cửa sổ Hanning Rw[tau]/Rw[0]
    \draw[thick, red!80!black, smooth] plot coordinates {
        (0,1.000) (15,0.996) (26,0.989) (40,0.973) (60,0.943) (80,0.902) 
        (100,0.851) (120,0.791) (140,0.725) (160,0.654) (180,0.584) (200,0.514) (214,0.468)
    };
    % Chú thích Rw[tau] đặt ở vùng thoáng (tau in [25, 65], y in [0.45, 0.75]) với mũi tên chỉ lên đường cong đỏ
    \node[red!80!black, font=\scriptsize, fill=white, inner sep=2.5pt, rounded corners=1pt, draw=red!40] (lbl_rw) at (46,0.60) 
        {$R_w[\tau] / R_w[0]$ (Cửa sổ co ngắn dần)};
    \draw[->, thick, red!80!black] (lbl_rw.north) -- (46,0.965);

    % Đường cong 2: Tự tương quan chưa chuẩn hóa Rxw[tau]/Rxw[0]
    \draw[thick, gray!60!black, dashed, smooth] plot coordinates {
        (0,1.000) (20,0.05) (40,-0.45) (60,-0.30) (80,0.25) (95,0.80) 
        (100,0.851) (105,0.80) (120,0.15) (140,-0.35) (160,-0.20) (180,0.15) 
        (195,0.48) (200,0.514) (205,0.48) (214,0.20)
    };
    % Chú thích Rxw đặt ở khoảng thoáng trung tâm (tau in [122, 154], y in [0.05, 0.19]), mũi tên chỉ thẳng xuống đáy đường cong xám
    \node[gray!70!black, font=\scriptsize, align=center, fill=white, inner sep=2.5pt, rounded corners=1pt, draw=gray!40] (lbl_rxw) at (138,0.12) 
        {$R_{xw}[\tau] / R_{xw}[0]$\\(Chưa chuẩn hóa: đỉnh bị dìm)};
    \draw[->, thick, gray!60!black] (lbl_rxw.south) -- (138,-0.33);

    % Đường cong 3: Sau khi chuẩn hóa U[tau] / R0n
    \draw[very thick, blue!85!black, smooth] plot coordinates {
        (0,1.000) (20,0.05) (40,-0.46) (60,-0.32) (80,0.28) (95,0.94) 
        (100,1.000) (105,0.94) (120,0.19) (140,-0.48) (160,-0.31) (180,0.26) 
        (195,0.93) (200,1.000) (205,0.93) (214,0.43)
    };
    \node[blue!85!black, font=\scriptsize, fill=white, inner sep=2pt, rounded corners=1pt] at (75,1.06) 
        {Đỉnh $\tau_0 = 100$ phục hồi $1{,}0$};
    \node[blue!85!black, font=\scriptsize, fill=white, inner sep=2pt, rounded corners=1pt] at (175,1.06) 
        {Đỉnh $2\tau_0 = 200$ phục hồi $1{,}0$};

    % Điểm mốc và mũi tên chỉ thị
    \filldraw[red!80!black] (214,0.468) circle (1.5pt);
    \node[above right, red!80!black, font=\scriptsize, fill=white, inner sep=1.5pt] at (214,0.50) 
        {$0{,}468$ tại $\tau=214$};

    \filldraw[blue!85!black] (100,1.0) circle (1.5pt);
    \filldraw[blue!85!black] (200,1.0) circle (1.5pt);
    \draw[->, thick, green!60!black] (100,0.851) -- (100,0.98);
    \draw[->, thick, green!60!black] (200,0.514) -- (200,0.98);

    \node[right, green!50!black, font=\scriptsize, fill=white, inner sep=1.5pt] at (102,0.92) {Bù $\times 1{,}175$};
    \node[left, green!50!black, font=\scriptsize, fill=white, inner sep=1.5pt] at (197,0.75) {Bù $\times 1{,}946$};
\end{tikzpicture}
}
\caption{Hiện tượng suy giảm tự nhiên của hàm tự tương quan cửa sổ $R_w[\tau]$ và sự khôi phục đỉnh chu kỳ cơ bản sau khi nhân bù hệ số $C[\tau]$ trên khung $N = 638\text{ mẫu}$ (số liệu trích xuất từ \texttt{verify_v3/ch08_norm_math.log}).}
\label{fig:ch08_window_effect}
\end{figure}

\textbf{Đọc hình.} Hình~\ref{fig:ch08_window_effect} minh họa trực quan ba đường cong tương quan trên dải trễ $\tau \in [0, 220]$. Đường cong màu đỏ liền nét biểu diễn tỉ số suy giảm hình học của cửa sổ Hanning $R_w[\tau]/R_w[0]$: đường cong dốc xuống liên tục từ $1{,}0$ tại $\tau = 0$ về mức $0{,}468$ tại $\tau = 214$. Đường cong nét đứt màu xám là hàm tự tương quan chưa chuẩn hóa $R_{xw}[\tau]/R_{xw}[0]$ của một nguyên âm có chu kỳ $\tau_0 = 100\text{ mẫu}$: đỉnh tại chu kỳ cơ bản $\tau_0 = 100$ bị dìm xuống còn $0{,}851$, và đỉnh tại chu kỳ bội $2\tau_0 = 200$ bị dìm nặng nề xuống mức $0{,}514$. Đường màu xanh lam đậm biểu thị kết quả sau khi chuẩn hóa cửa sổ: đỉnh tại $\tau_0 = 100$ được nhân bù $\times 1{,}175$ và phục hồi về đúng $1{,}0$; đỉnh tại $2\tau_0 = 200$ được nhân bù $\times 1{,}946$ và cũng phục hồi về mức $1{,}0$. Dải quét phần cứng $\tau \in [26, 214]$ được đóng khung nền xanh nhạt, bao trọn toàn bộ miền tần số giọng nói $75\text{ Hz}$ đến $600\text{ Hz}$.

\textit{Tóm lại: Cắt khung bằng cửa sổ Hanning làm suy giảm hình học hàm tự tương quan khi độ trễ tăng; nếu không chia bù cho $R_w[\tau]$, các đỉnh ở chu kỳ dài bị dìm sâu và gây lỗi nhảy quãng tám lên trên; việc nhân với nghịch đảo $C[\tau] = R_w[0]/R_w[\tau]$ khôi phục trung thực độ mạnh của các đỉnh chu kỳ cơ bản.}
